\section{National trajectory and source-support specifications}\label{sec:SI_national}
Actual national trajectories use all 183 eligible identifiers in 2015--2023; the separate 2010--2020 account comparison retains 182 eligible members of that cohort. Account contributions are $100\Delta B_k/B_{\rm total,0}$ percentage points. Forest-component decomposition uses the symmetric product identity for source account area and the implied gha/ha factor. That factor belongs to the accounting lineage and is not independent forest productivity or stock evidence. The account-domain and component contribution tables retain all observed signs.

National MapBiomas tables are summed at non-overlapping leaf classes within each published geography. Brazilian formations are $\{3,5,6,49\}$, with $\{4,7\}$ added only for the distinct formations-plus-savanna analysis. Indonesian forest is $\{3,5,76\}$, oil palm is $\{35\}$ and the broader agricultural target is $\{9,21,35,40\}$. Congolese source forest is $\{3,4,5\}$, dense forest plus mangrove is $\{3,5\}$ and plantation remains $\{9\}$. These are not interchangeable forest definitions. All source footprints, zero/no-data conventions and class crosswalks remain in the data.

The Brazilian reported footprint is 850,644,939.994 ha and the Indonesian Collection 4.1 footprint is 189,044,390.245 ha. Consecutive transition margins reproduce finest-geography extent totals to less than 0.000005 and 0.0000003 ha, respectively. Those residuals measure arithmetic consistency, not map accuracy. Indonesia's Collection 3 footprint exceeds Collection 4.1 by 43,130.051 ha. A difference of total areas is not a geometry symmetric difference, and source/support differences cannot be isolated as a pure collection effect.

The Congo archive folder is named ``wide'', but its CSV members have long schema with explicit year, class and area fields. The country member contains 234 rows, nine observed leaf classes and all 26 years, 2000--2025. Units are established by the version-pinned source script's \texttt{pixelArea().divide(10000)}, not inferred from an unlabelled \texttt{area} header. Source statistics round class areas to two decimal digits. The country label footprint ranges from 231,404,518.04 to 231,404,518.08 ha across years. Biome and province exports do not sum to this footprint; they remain separate support summaries. Neither near-constant label area nor net footprint differences certify a complete, disjoint geometry partition or the extent of unobserved forest.

Congo's country forest-group endpoint difference is $-3,303,015.650$ ha ($-2.014673\%$). The symmetric endpoint net-area break-even requirement is $1,651,507.825$ ha per endpoint, not an uncertainty estimate. Its ATBD describes a 60-km corridor on each side of the Miombo--Mosaic boundary, additional classification trained on temporally stable map labels outside the corridor, and replacement restricted to selected vegetation classes. It also documents class confusion and irregular temporal series. These procedures support the provenance description, not independent paired change validation. No stock loss or forest-to-agriculture conversion is inferred from the marginal area changes.

\begin{figure}[htbp]
\centering\includegraphics[width=1\textwidth]{Indonesia_definition_support_SI.png}
\caption{Indonesia definition and collection/support sensitivity, 2015--2023. Left: changes in source forest subclasses and their union. Right: C3 versus C4.1 mapped forest changes; differing support prevents interpretation as a pure collection effect or confidence interval. Main national results summarize this sensitivity. Source: MapBiomas Indonesia; original numerical figure retained.}\label{fig:SI_Indonesia}
\end{figure}
\begin{table}[htbp]
\centering\small
\caption{Source-specific class groups and geographic support checks. Numeric class codes are initiative-specific. Footprint differences are not geometry symmetric differences.}\label{tab:SI_class_support}
\begin{tabularx}{\textwidth}{lp{4.0cm}X}\toprule
Source & Mapped group & Support qualification\\\midrule
Brazil C11 & 3, 5, 6, 49; plus 4, 7 only for F+S & 850.645 million ha; annual margins reconciled.\\
Indonesia C4.1 & 3, 5, 76; oil palm 35 separate & 189.044 million ha; C3 footprint 43,130.051 ha larger.\\
DR Congo C1 & 3, 4, 5; plantation 9 separate & Country footprint 231.405 million ha; province/biome totals not forced into a partition.\\
\bottomrule\end{tabularx}
\end{table}

\begin{table}[htbp]
\centering\small
\caption{Exact Indonesian mapped class and collection/support sensitivities, 2015--2023. Areas at each endpoint are million ha; changes are ha and percentage of the initial class area. C3 and C4.1 do not have certified identical support, so their contrast is not a pure version effect or an uncertainty interval.}\label{tab:SI_Indonesia_classes}
\begin{tabularx}{\textwidth}{lrrrr}\toprule
Source/class & 2015 (Mha) & 2023 (Mha) & Change (ha) & Change (\%)\\\midrule
C4.1 Forest Formation & 93.124119 & 93.127138 & +3019.087 & +0.0032\\
C4.1 mangrove & 3.007685 & 2.996065 & -11620.148 & -0.3863\\
C4.1 peat swamp forest & 10.922995 & 10.229033 & -693961.333 & -6.3532\\
C4.1 combined forest & 107.054799 & 106.352236 & -702562.394 & -0.6563\\
C4.1 oil palm (separate) & 14.880928 & 17.714296 & +2833367.219 & +19.0403\\
C3 level-one Forest & 98.031838 & 94.925798 & -3106039.930 & -3.1684\\
\bottomrule\end{tabularx}
\end{table}
\begin{table}[htbp]
\centering\small
\caption{Annual mapped transition-event sums, 2015--2023 (million hectare-events). Repetition across years and intermediate paths is allowed. These are not unique endpoint-cohort areas, calibrated conversion or causal effects.}\label{tab:SI_transition_events}
\begin{tabularx}{\textwidth}{lrrr}\toprule
Group/event & Outflow & Inflow & Inflow $-$ outflow\\\midrule
Brazil formations (all destinations) & 22.710844 & 11.848045 & -10.862799\\
Indonesia forest (all destinations) & 4.424120 & 3.721558 & -0.702562\\
Indonesia forest, direct to oil palm & 0.707004 & -- & --\\
\bottomrule\end{tabularx}
\end{table}
\begin{table}[htbp]
\centering\small
\caption{Mapped subnational net-change compensation, 2015--2023 (million ha). Negative and positive sums are over states/provinces, not annual gross transition events. They are descriptive classifications, not independent observations or confidence-resolved ecological gains and losses.}\label{tab:SI_regional_compensation}
\begin{tabularx}{\textwidth}{llrrrrr}\toprule
ISO & Partition & Declining & Increasing & Negative sum & Positive sum & National net\\\midrule
BRA & state (27) & 20 & 7 & -11.017319 & +0.154520 & -10.862799\\
IDN & province (38) & 14 & 24 & -1.303387 & +0.600825 & -0.702562\\
\bottomrule\end{tabularx}
\end{table}

\section{Sharp endpoint-cohort bounds and independent numerical certificates}
Let $G$ be a layered directed acyclic graph with vertices $(t,i)$ and edge capacities $T_{tij}\ge0$. For a marked cohort whose complete origin is the specified 2015 group, the feasible set is
\[
\mathcal X=\{x:0\le x\le T,\ x_{2015,ij}=T_{2015,ij}\mathbf1\{i\in F\},\ Ax=0\},
\]
where $A$ imposes conservation at every interior vertex. For a destination $D$, the objective is $c^Tx=\sum_i\sum_{j\in D}x_{2022,ij}$. The minimum and maximum give unique endpoint-cohort bounds, not event totals.

If capacities are exactly balanced, the unmarked flow $T-x$ is also nonnegative and conserved. Both flows decompose into paths; their continuum of area parcels constitutes a compatible land history. Conversely, any compatible history generates a feasible marked flow. Thus the bounds are sharp for continuous-area aggregate constraints, not necessarily for an integer original-pixel lattice or ecologically restricted histories. Any path may traverse intermediate classes or repeat conversion. Restricting flows to reported fixed geographic units removes cross-unit reallocations. Its projected feasible set is a subset of the national aggregate set, so raw extrema cannot weaken. Numerical outward allowances may make rounded endpoints marginally wider and are reported separately.

The seven group--target cases comprise four Brazilian and three Indonesian cases. Two optimization directions at national level and within each of 49 or 38 units yield 634 programmes. Hectares are scaled by $10^5$ inside the solver. The recorded feasibility and gap limit is 0.02 ha per programme, including $L^1$ conservation, source imbalance and bound-violation diagnostics. The independent witness check reconstructs constraints and verifies primal feasibility, dual stationarity, complementarity and primal--dual objective agreement from the deposited capacities without invoking the optimizer. Source-balanced ideal-model sharpness and finite-precision numerical certificates are distinct claims.

The finest-geography endpoint bounds in million ha are:
\begin{table}[htbp]
\centering\small
\caption{Exact fine-geography endpoint-cohort ranges (million ha), corresponding to main Figure 5. Values include outward numerical allowances; they are feasible mapped ranges, not confidence intervals.}\label{tab:SI_conversion_exact}
\begin{tabularx}{\textwidth}{lllrr}\toprule
Country & Initial group & 2023 target & Lower & Upper\\\midrule
BRA & formations & nonforest & 11.203746901 & 22.710751994\\
BRA & formations & cropland & 0.000000000 & 7.014430883\\
BRA & formations plus savanna & nonforest & 18.373845177 & 35.847255651\\
BRA & formations plus savanna & cropland & 0.009385753 & 8.661678113\\
IDN & mapped forest & nonforest & 1.312001539 & 4.424121169\\
IDN & mapped forest & oil palm & 0.111612129 & 2.325782972\\
IDN & mapped forest & agriculture & 0.126800213 & 3.763620519\\
\bottomrule\end{tabularx}
\end{table}


For true endpoint sets $F,D$ and mapped sets $\widehat F,\widehat D$, a union-set argument gives
\[
\bigl|\mu(F\cap D)-\mu(\widehat F\cap\widehat D)\bigr|
\le\mu(F\mathbin{\triangle}\widehat F)+\mu(D\mathbin{\triangle}\widehat D).
\]
Consequently a mapped lower bound $L$ transfers only as $\max(0,L-E)$ for a compatible combined set-error budget $E$. Indonesia's oil-palm lower bound requires $E<111,612.129$ ha; an equal split would be approximately 55,806 ha per endpoint. Neither an overall classification score nor net endpoint area error establishes this requirement.

\begin{table}[htbp]
\centering\small
\caption{Finnish inventory source cells and reported volume standard errors. Units are million m$^3$ after conversion from the source thousand-m$^3$ cells; source identities are pinned in the numerical dependency.}\label{tab:SI_inventory_cells}
\begin{tabularx}{\textwidth}{lrrll}\toprule
Inventory & Volume & SE & Sheet & Cells\\\midrule
NFI12 & 2475.192 & 13.001 & 20a & N61/O61\\
NFI13 & 2551.964 & 13.534 & Table 55 & N60/O60\\
\bottomrule\end{tabularx}
\end{table}

\section{Inventory source cells, period operators and dependence bounds}
Finnish original source cells are NFI12 supplementary \texttt{20a!N61/O61}, 2,475,192 and 13,001 thousand m$^3$, and NFI13 \texttt{Table55!N60/O60}, 2,551,964 and 13,534 thousand m$^3$. The values after the slash are sampling standard errors, not 95\% halfwidths. Both inventories include \AA land. The complete source table identifies units and inventory cycles. The approximate normal Bonferroni difference band, including source rounding, is $[17.293137220,136.250862780]$ million m$^3$. A differential-bias budget of 17.293137220 million m$^3$ would exhaust its positive lower bound. Separately, $\Delta\widehat V\pm\sqrt{20}(s_0+s_1)$ is a worst-covariance second-moment sensitivity: the worst-variance standard deviation is bounded by $s_0+s_1$, and the $\sqrt{20}$ multiplier is the nominal 95\% Chebyshev form. It gives $[-41.896127566,195.440127566]$ million m$^3$, conditional on unbiased errors and reported standard errors equal to true error standard deviations. Exact coverage is not asserted for estimated variances. Both constructions exclude unbounded differential bias.

The Finnish biomass API reports 1,744 and 1,786 Mt for total living trees including below-ground components. Independently rounded stems plus branches/foliage sum to 1,350 and 1,384 Mt; independently rounded components can differ from the published total by approximately one million tonnes. No biomass standard error is inferred from volume error. 

French 2013--2017 volume and 95\% halfwidth are 2,716 and 45 million m$^3$ in the France row of \emph{M\'emento 2018}, printed pp.~28--29, PDF page 16 (a two-page spread). The row is the national France row, not \^Ile-de-France. The associated methodology, PDF pp.~5, 9, 16--17, describes five-campaign pooling, protocol changes, production-forest categories including poplar, the diameter threshold and unusable-wood deduction. Current workbook national totals are read from \texttt{BV Propri\'et\'e}, row \texttt{Total}, in the respective Volume XLSX member:
\begin{table}[htbp]\centering\small
\caption{French published national volume estimates and 95\% marginal halfwidths. Campaign windows are pooled observations, not exact annual endpoints.}\label{tab:SI_France_windows}
\begin{tabular}{lrr}\toprule
Campaign window & Estimate (million m$^3$) & Published 95\% halfwidth\\\midrule
2013--2017 & 2,716 & 45\\
2018--2022 & 2,811 & 46\\
2019--2023 & 2,827 & 46\\
2021--2025 & 2,916 & 47\\\bottomrule
\end{tabular}\end{table}

French workbook sheet dimensions incorrectly advertise \texttt{A1:A1}; parsing resets those dimensions before iterating the actual seven columns. The readme specifies 95\% intervals, production forest and the stock definition. The latest includes 30,136 inventory points, 29,502 outside poplar plantations and 634 in poplar plantations. Changed physical-accessibility and planted-forest breakdowns are not compared across old/new classifications. Matching nominal stock definitions does not certify exact retrospective cubage or unusable-wood harmonization. A residual differential-method bias budget is left explicit rather than set to zero.

With full-unit allowance $r=1$ million m$^3$ for each estimate and each halfwidth, the per-contrast 90\% band is
\[
\Delta\widehat V\pm[h_0+h_1+4r].
\]
For an approximate 95\% marginal-normal band, use
\[
\Delta\widehat V\pm\left[\frac{z_{1-0.05/(2m)}}{z_{0.975}}(h_0+h_1+2r)+2r\right],
\]
with $m=2$ for one contrast or $m=4$ for simultaneous coverage of the four endpoints and hence all three latest-cycle contrasts. No independence is assumed. The longest difference is $+200$ million m$^3$: the per-contrast 90\% band is $[104,296]$, the pairwise approximate 95\% band is $[90.502180,309.497820]$, and the four-endpoint approximate 95\% band is $[78.209883,321.790117]$. All three earlier-to-latest contrasts remain in the deposited table, including the shorter unresolved ones. A systematic differential bias with absolute bound $B$ expands either side by $B$; no measured bound is asserted. These are conditional inventory-volume results, not calibrated national 2015--2023 AGB.

\begin{table}[htbp]
\centering\small
\caption{Finnish forest-account timber-stock reconciliation (million m$^3$). The final displayed component combines reclassification and balancing. In 2023, the opening-to-closing change is $-4.185$ despite net increment minus removals/losses of $+20.802$, because that component is $-24.987$. This partly inventory-derived account is neither independent biomass validation nor AGB mass.}\label{tab:SI_FIN_stock_flow}
\begin{tabularx}{\textwidth}{lrrrrr}\toprule
Year & Opening & Closing & Stock change & Increment $-$ removals & Reclass./balancing\\\midrule
2022 & 2588.158 & 2642.554 & +54.396 & +20.808 & +33.588\\
2023 & 2642.553 & 2638.368 & -4.185 & +20.802 & -24.987\\
\bottomrule\end{tabularx}
\end{table}

\section{Canadian biomass change and uncertainty representations}
The Canadian source is NFI version 3.1, November 2025. Tables 17.0 and 17C.0 report stock and annualized change by forest type and age class; accompanying \texttt{sep} and \texttt{intv} files report relative standard errors and interval halfwidths. CSV values retain greater numerical precision than their HTML renderings. The HTML footers establish combined remeasured 2007--2017 and projected 2006--2017 data, removal of strong artificial-change indicators, and omitted Arctic ecozones plus the Nunavut portions of Taiga Plains and Hudson Plains. The omitted area is 262,189.89 thousand ha, not assumed to carry zero forest biomass.

The stock total is 30,525.199 Mt with labelled RSE 0.70\% and halfwidth 418.4005 Mt, approximately consistent with the usual large-sample interval relation after RSE rounding. Annualized total change is $+17.639$ Mt/year with labelled RSE 1.16\% and halfwidth 38.9180 Mt/year. The latter ratio $38.918/(17.639\times0.0116)$ is 190.204, not approximately 2. Individual change rows also fail the ordinary relation, so a rounded percentage alone cannot supply a change standard error. The point, interval and RSE HTML tables all carry the same version 3.1, November 2025 label; a simple labelled-version mismatch does not resolve the discrepancy. The published halfwidth is retained as labelled and its confidence level is not independently certified for this export. The general estimation document uses default Type-I probability 0.05 and a national critical value near 2, but it is not a current implementation specification for these projected change tables.

\begin{table}[htbp]\centering\small
\caption{Canadian annualized change by contemporaneous forest type, with provider-labelled halfwidths. No confidence level is assigned here; components are not fixed biological cohorts.}\label{tab:SI_Canada_components}
\begin{tabular}{lrr}\toprule
Source forest type & Annualized change (Mt/year) & Labelled halfwidth\\\midrule
Coniferous & $-16.158$ & 34.1130\\
Mixedwood & $-16.115$ & 12.2560\\
Broadleaf & $+43.065$ & 13.9775\\
Non-treed & $+6.847$ & 2.3140\\
Total & $+17.639$ & 38.9180\\\bottomrule
\end{tabular}\end{table}

All four component rows and the total are retained; none is selected as a confirmatory sign test. They refer to stocks classified by contemporaneous forest type, not fixed biological cohorts. Source totals reconcile to rounding. The direct total interval $[-21.279,56.557]$ Mt/year includes zero, and full-CSV-unit outward rounding gives $[-21.2801,56.5581]$. Projection, classification error and artificial-change filtering remain separate from sampling precision. The source design permits useful inventory-domain evidence without imposing an exact paired census, but it does not support transfer to omitted geography or the primary endpoints.

\section{Reference compatibility protocols}
The Brazilian Collection 11 accuracy appendix uses reference collection 3 for 1985--2024, with boundary-point exclusion, interpreter weighting, sparse-class exclusions and map-dependent legend harmonization. The examined public collection-2.1 archive is an older reference representation. Of 85,152 points, 4,053 have nonblank 2015/2023 pairs; 4,041 remain after removing non-observed labels and the archived training role; 3,042 are also nonedge at both endpoints. These counts do not certify a compatible national probability-design sample. Local zero paired-reference eligibility remains a property of the original footprint and examined artifacts, not a global absence claim. The different current/old training counts are preserved as provenance differences, not reconciled by assumption. A newly retrieved vector archive currently linked from the official download page has an October 2026 URL but 85,152 records and 126 fields ending in 2022; no 2023 or 2024 annual fields occur. Its recent URL therefore does not establish the current reference collection or fill the endpoint gap. The DBF header is checked without expanding the full table.

Indonesia's Collection 4.1 ATBD describes independent annual references through 2022 and does not supply a current global Collection 4.1 accuracy estimate in its collection summary. Older pooled accuracy does not calibrate its 2015--2023 change. Congo's ATBD describes limitations and future validation improvements but does not establish a compatible endpoint change interval. None of these findings proves that an adequate reference design is impossible. Compatible marginal reference intervals, direct inventory-change estimates, defensible covariance bounds or bounded bias can be sufficient; exact paired field censuses and known covariance are not universal prerequisites.

\begin{table}[htbp]
\centering\small
\caption{All French earlier-to-2021--2025 volume contrasts (million m$^3$), including full-unit rounding. Normal rows are approximate and conditional on marginal normality. An unbounded differential-method bias is not covered.}\label{tab:SI_France_contrasts}
\begin{tabularx}{\textwidth}{llrrr}\toprule
Initial cycle & Construction & Difference & Lower & Upper\\\midrule
2013--2017 & 90\% union & 200 & 104.000000 & 296.000000\\
2013--2017 & 95\% normal, 2 ends & 200 & 90.502180 & 309.497820\\
2013--2017 & 95\% normal, 4 ends & 200 & 78.209883 & 321.790117\\
2018--2022 & 90\% union & 105 & 8.000000 & 202.000000\\
2018--2022 & 95\% normal, 2 ends & 105 & -5.641414 & 215.641414\\
2018--2022 & 95\% normal, 4 ends & 105 & -18.064480 & 228.064480\\
2019--2023 & 90\% union & 89 & -8.000000 & 186.000000\\
2019--2023 & 95\% normal, 2 ends & 89 & -21.641414 & 199.641414\\
2019--2023 & 95\% normal, 4 ends & 89 & -34.064480 & 212.064480\\
\bottomrule\end{tabularx}
\end{table}
